\begin{appendix}
\section*{附录A FILM网络核心代码}

本附录提供了FILM网络模型的核心实现代码，便于读者理解和复现本文所采用的方法。

\subsection*{A.1 跨注意力融合块（restormer\_cablock）}

\begin{lstlisting}[language=Python]
class restormer_cablock(nn.Module):
    def __init__(self, input_channel=1, restormerdim=32,
                 restormerhead=8, image2text_dim=10,
                 ffn_expansion_factor=4, bias=False,
                 LayerNorm_type='WithBias', hidden_dim=768,
                 pooling='avg', normalization='l1'):
        super().__init__()
        self.convA1 = nn.Conv2d(input_channel, restormerdim,
                                kernel_size=3, stride=1, padding=1)
        self.preluA1 = nn.PReLU()
        self.convA2 = nn.Conv2d(image2text_dim, restormerdim, kernel_size=1)
        self.preluA2 = nn.PReLU()
        self.convA3 = nn.Conv2d(2 * restormerdim, restormerdim, kernel_size=1)
        self.preluA3 = nn.PReLU()
        # 图像B分支（对称结构）...
        self.restormerA1 = Restormer(restormerdim, restormerhead,
                                     ffn_expansion_factor, bias, LayerNorm_type)
        self.cross_attentionA1 = CrossAttention(embed_dim=hidden_dim, num_heads=8)
        self.imagef2textfA1 = imagefeature2textfeature(restormerdim, image2text_dim, hidden_dim)

    def forward(self, imageA, imageB, text):
        # Restormer特征提取
        imageA = self.restormerA1(self.preluA1(self.convA1(imageA)))
        imageAtotext = self.imagef2textfA1(imageA)
        # 跨注意力计算
        ca_A = self.cross_attentionA1(text, imageAtotext, imageAtotext)
        # 注意力池化与归一化
        ca_A = F.adaptive_avg_pool1d(ca_A.permute(0, 2, 1), 1).permute(0, 2, 1)
        ca_A = F.normalize(ca_A, p=1, dim=2)
        # 特征重构与融合
        ca_A = (imageAtotext * ca_A).view(b, self.image2text_dim, 288, 384)
        ca_A = F.interpolate(ca_A, [H, W], mode='nearest')
        ca_A = self.preluA3(self.convA3(torch.cat(
            (F.interpolate(imageA, [H, W]),
             self.preluA2(self.convA2(ca_A)) + imageA_sideout), 1)))
        return ca_A, ca_B
\end{lstlisting}

\subsection*{A.2 FILM主网络（Net）}

\begin{lstlisting}[language=Python]
class Net(nn.Module):
    def __init__(self, mid_channel=32, decoder_num_heads=8,
                 ffn_factor=4, bias=False,
                 LayerNorm_type='WithBias', out_channel=1,
                 hidden_dim=256, image2text_dim=32):
        super().__init__()
        self.text_process = text_preprocess(768, hidden_dim)
        self.restormerca1 = restormer_cablock(hidden_dim=hidden_dim,
                                               image2text_dim=image2text_dim)
        self.restormerca2 = restormer_cablock(input_channel=mid_channel,
                                               hidden_dim=hidden_dim,
                                               image2text_dim=image2text_dim)
        self.restormerca3 = restormer_cablock(input_channel=mid_channel,
                                               hidden_dim=hidden_dim,
                                               image2text_dim=image2text_dim)
        self.restormer1 = Restormer(2 * mid_channel, decoder_num_heads,
                                    ffn_factor, bias, LayerNorm_type)
        self.restormer2 = Restormer(mid_channel, decoder_num_heads,
                                    ffn_factor, bias, LayerNorm_type)
        self.restormer3 = Restormer(mid_channel, decoder_num_heads,
                                    ffn_factor, bias, LayerNorm_type)
        self.conv1 = nn.Conv2d(2 * mid_channel, mid_channel, kernel_size=1)
        self.conv2 = nn.Conv2d(mid_channel, out_channel, kernel_size=1)
        self.softmax = nn.Sigmoid()

    def forward(self, imageA, imageB, text):
        text = self.text_process(text)
        featureA, featureB = self.restormerca1(imageA, imageB, text)
        featureA, featureB = self.restormerca2(featureA, featureB, text)
        featureA, featureB = self.restormerca3(featureA, featureB, text)
        fusionfeature = torch.cat((featureA, featureB), 1)
        fusionfeature = self.restormer1(fusionfeature)
        fusionfeature = self.conv1(fusionfeature)
        fusionfeature = self.restormer2(fusionfeature)
        fusionfeature = self.restormer3(fusionfeature)
        fusionfeature = self.conv2(fusionfeature)
        fusionfeature = self.softmax(fusionfeature)
        return fusionfeature
\end{lstlisting}

\section*{附录B 系统API接口说明}

\subsection*{B.1 AI融合接口}

\textbf{POST /api/fuse}

请求参数：FormData格式，包含over\_image和under\_image两个文件字段。

响应：image/jpeg格式的融合结果图像。

错误码：400（文件类型/大小/尺寸不匹配）、500（服务器内部错误）。

\subsection*{B.2 传统融合接口}

\textbf{POST /api/fuse/traditional}

请求参数：FormData格式，包含over\_image、under\_image和algorithm三个字段，algorithm可选值为avg、max、mertens。

响应：image/jpeg格式的融合结果图像。

\subsection*{B.3 评估接口}

\textbf{POST /api/evaluate}

请求参数：FormData格式，包含fused\_image文件字段。

响应：JSON格式，包含en、sd、sf、ag四项指标。

\subsection*{B.4 历史查询接口}

\textbf{GET /api/history}

请求参数：无。

响应：JSON数组，包含最近20条历史记录。
\end{appendix}
